\section{Eigenvalues and Eigenvectors}

\subsection{Invariant Subspaces}

\begin{exercise}[5a1]
    Suppose $T \in \lin{V}$ and $U$ is a subspace of $V$.
    \begin{subexercises}
        \item Prove that if $U \subseteq \nul T$, then $U$ is invariant under $T$.
        \item Prove that if $\range T \subseteq U$, then $U$ is invariant under $T$.
    \end{subexercises}
\end{exercise}

\begin{sole}
    \item If $U \subseteq \nul T$, then for any $u \in U$, we have $T(u) = 0 \in U$, so $U$ is invariant under $T$.
    \item If $\range T \subseteq U$, then for any $u \in U$, $T(u) \in \range T \subseteq U$, so $U$ is invariant under $T$.
\end{sole}

\begin{exercise}[5a2]
    Suppose that $T \in \lin{V}$ and $V_1, \ldots, V_m$ are subspaces of $V$ invariant under $T$.
	Prove that $V_1 + \cdots + V_m$ is invariant under $T$.
\end{exercise}

\begin{sol}
    Let $v \in V_1 + \cdots + V_m$.
	Then $v = v_1 + \cdots + v_m$ for some $v_i \in V_i$.
	Since each $V_i$ is invariant under $T$, we have $T(v_i) \in V_i$ for each $i$.
	Therefore,
    \[T(v) = T(v_1 + \cdots + v_m) = T(v_1) + \cdots + T(v_m) \in V_1 + \cdots + V_m,\]
    so $V_1 + \cdots + V_m$ is invariant under $T$.
\end{sol}

\begin{exercise}[5a3]
    Suppose $T \in \lin{V}$.
	Prove that the intersection of every collection of subspaces of $V$ invariant under $T$ is invariant under $T$.
\end{exercise}

\begin{sol}
    Let $\{U_i\}_{i \in I}$ be a collection of subspaces of $V$ invariant under $T$, and let
    \[U = \bigcap_{i \in I} U_i.\]
    Let $u \in U$.
	Then $u \in U_i$ for all $i \in I$.
	Since each $U_i$ is invariant under $T$, we have $T(u) \in U_i$ for all $i \in I$.
	Therefore, $T(u) \in U$, so $U$ is invariant under $T$.
\end{sol}

\begin{exercise}[5a4]
    Prove or give a counterexample: If $V$ is finite-dimensional and $U$ is a subspace of $V$ that is invariant under every operator on $V$, then $U = \set{0}$ or $U = V$.
\end{exercise}

\begin{sol}
    If $U = \set 0$, then we are done.
	Now suppose $U \neq \set 0$ so that there exists nonzero $v_1 \in U$.
	Extend $v_1$ to a basis $v_1, \ldots, v_m$ of $V$, and let $v \in V$.
	Define the linear operator $T \in \lin{V}$ by $T(v_1) = v$ and $T(v_i) = 0$ for $i = 2, \ldots, m$.
	Since $U$ is invariant under $T$, we have $T(v_1) = v \in U$.
	But $v$ was an arbitrary element of $V$, so $U = V$.
	Therefore, $U = \set{0}$ or $U = V$.
\end{sol}

\begin{exercise}[5a5]
    Suppose $T \in \lin{\rbb^2}$ is defined by $T(x, y) = (-3y, x)$.
	Find the eigenvalues of $T$.
\end{exercise}

\begin{sol}
    If $(x, y)$ is an eigenvector of $T$ with eigenvalue $\lambda$, then we have the system of equations
    \[
    \begin{cases}
        -3y = \lambda x,\\
        x = \lambda y
    \end{cases}
    \]
    Solving the system, we get $\lambda^2 = -3$ which has no real solutions.
	Therefore, $T$ has no real eigenvalues.
\end{sol}

\begin{exercise}[5a6]
    Define $T \in \lin{\fbb^2}$ by $T(w, z) = (z, w)$.
	Find all eigenvalues and eigenvectors of $T$.
\end{exercise}

\begin{sol}
    If $(w, z)$ is an eigenvector of $T$ with eigenvalue $\lambda$, then we have
    \[
    (z, w) = \lambda (w, z),
    \]
    which gives the system
    \[
    \begin{cases}
        z = \lambda w,\\
        w = \lambda z.
    \end{cases}
    \]
    Solving, we get $\lambda^2 = 1$, so $\lambda = 1$ or $\lambda = -1$.
	If $\lambda = 1$, then $z = w$, so the eigenvectors are of the form $(w, w)$.
	If $\lambda = -1$, then $z = -w$, so the eigenvectors are of the form $(w, -w)$.
\end{sol}

\begin{exercise}[5a7]
    Define $T \in \lin{\fbb^3}$ by $T(z_1, z_2, z_3) = (2z_2, 0, 5z_3)$.
	Find all eigenvalues and eigenvectors of $T$.
\end{exercise}

\begin{sol}
    If $(z_1, z_2, z_3)$ is an eigenvector of $T$ with eigenvalue $\lambda$, then we have
    \[
    (2z_2, 0, 5z_3) = \lambda (z_1, z_2, z_3),
    \]
    which gives the system
    \[
    \begin{cases}
        2z_2 = \lambda z_1,\\
        0 = \lambda z_2,\\
        5z_3 = \lambda z_3.
    \end{cases}
    \]
    Solving, we get the eigenvalues $\lambda = 0$ and $\lambda = 5$.
	For $\lambda = 0$, the eigenvectors are of the form $(z_1, 0, 0)$.
	For $\lambda = 5$, the eigenvectors are of the form $(0, 0, z_3)$.
\end{sol}

\begin{exercise}[5a8]
    Suppose $P \in \lin{V}$ is such that $P^2 = P$.
	Prove that if $\lambda$ is an eigenvalue of $P$, then $\lambda = 0$ or $\lambda = 1$.
\end{exercise}

\begin{sol}
    If $v$ is an eigenvector of $P$ with eigenvalue $\lambda$, then $Pv = \lambda v$.
	Applying $P$ again, we get $P^2v = P(\lambda v) = \lambda Pv = \lambda^2 v$.
	But $P^2 = P$, so $P^2v = Pv = \lambda v$.
	Therefore, $\lambda^2 v = \lambda v$, which gives $(\lambda^2 - \lambda)v = 0$.
	Since $v \neq 0$, we must have $\lambda^2 - \lambda = 0$, so $\lambda(\lambda - 1) = 0$.
	Hence, $\lambda = 0$ or $\lambda = 1$.
\end{sol}

\begin{exercise}[5a9]
    Define $T: \pcal(\rbb) \to \pcal(\rbb)$ by $Tp = p'$.
	Find all eigenvalues and eigenvectors of $T$.
\end{exercise}

\begin{sol}
    Let $p \in \pcal(\rbb)$ such that $p' = \lambda p$ for some $\lambda \in \rbb$.
    \begin{itemize}
        \item If $\lambda \neq 0$, suppose $\deg p = 0$. Then $p$ is a constant polynomial, and $p' = 0 \neq \lambda p$, so $\deg p \geq 1$.
        If $\deg p = n \geq 1$, then $\deg p' = n - 1$, so $\deg p' \neq \deg p$, which contradicts $p' = \lambda p$.
        \item If $\lambda = 0$, then $p' = 0$, which implies that $p$ is a constant polynomial.
	Therefore, the eigenvectors corresponding to the eigenvalue $\lambda = 0$ are all nonzero constant polynomials. \qh
    \end{itemize}
\end{sol}

\begin{exercise}[5a10]
    Define $T \in \lin{\pcal_4(\rbb)}$ by $(Tp)(x) = xp'(x)$ for all $x \in \rbb$.
	Find all eigenvalues and eigenvectors of $T$.
\end{exercise}

\begin{sol}
    Let $p(x) = a_0 + a_1x + a_2x^2 + a_3x^3 + a_4x^4$ be an eigenvector of $T$ with eigenvalue $\lambda$.
    Then
    \[(Tp)(x) = xp'(x) = x(a_1 + 2a_2x + 3a_3x^2 + 4a_4x^3) = a_1x + 2a_2x^2 + 3a_3x^3 + 4a_4x^4.\]
    Equating coefficients with $\lambda p(x) = \lambda a_0 + \lambda a_1x + \lambda a_2x^2 + \lambda a_3x^3 + \lambda a_4x^4$ yields the system:
    \[\lambda a_0 = 0, \quad \lambda a_1 = a_1, \quad \lambda a_2 = 2a_2, \quad \lambda a_3 = 3a_3, \quad \lambda a_4 = 4a_4.\]
    For each $n \in \set{0, 1, 2, 3, 4}$, the equation $\lambda a_n = n a_n$ is satisfied with $a_n \neq 0$ if and only if $\lambda = n$.
    When $\lambda = n$, all other coefficients must be zero, so the eigenvector is a nonzero scalar multiple of $x^n$.
    Therefore, the eigenvalues of $T$ are $0, 1, 2, 3, 4$, with corresponding eigenvectors being nonzero scalar multiples of $1, x, x^2, x^3, x^4$, respectively.
\end{sol}

\begin{exercise}[5a11]
    Suppose $V$ is finite-dimensional, $T \in \lin{V}$, and $\alpha \in \fbb$.
	Prove that there exists $\delta > 0$ such that $T - \lambda I$ is invertible for all $\lambda \in \fbb$ such that $0 < |\alpha - \lambda| < \delta$.
\end{exercise}

\begin{sol}
    Since $V$ is finite-dimensional, the set of eigenvalues of $T$ is finite.
	If $T$ has no eigenvalues, then $T - \lambda I$ is invertible for all $\lambda \in \fbb$, and we can choose any $\delta > 0$.
	Otherwise, let $\lambda_1, \ldots, \lambda_m$ be the eigenvalues of $T$.
	Define
    \[\delta = \min_{\lambda_j \neq \alpha} |\alpha - \lambda_j|.\]
    If $\alpha$ is the only eigenvalue, then we can select any $\delta > 0$ since any $\lambda$ satisfying $0 < |\alpha - \lambda| < \delta$ will not be equal to $\alpha$, and hence $T - \lambda I$ will be invertible.
	Otherwise, we have for any $\lambda \in \fbb$ such that $0 < |\alpha - \lambda| < \delta$, we have $\lambda \neq \lambda_j$ for all $j = 1, \ldots, m$, so $\lambda$ is not an eigenvalue of $T$.
	Hence, $T - \lambda I$ is invertible for all such $\lambda$.
\end{sol}

\begin{exercise}[5a12]
    Suppose $V = U \op W$, where $U$ and $W$ are nonzero subspaces of $V$.
	Define $P \in \lin{V}$ by $P(u + w) = u$ for each $u \in U$ and each $w \in W$.
	Find all eigenvalues and eigenvectors of $P$.
\end{exercise}

\begin{sol}
    Let $v = u + w \in V$ be an eigenvector of $P$ with eigenvalue $\lambda$.
	Then $P(v) = P(u + w) = u = \lambda v = \lambda (u + w)$.
	Comparing the $U$ and $W$ components, we get $u = \lambda u$ and $0 = \lambda w$.
	If $\lambda \neq 0$, then $w = 0$ and $u = \lambda u$ implies $\lambda = 1$.
	If $\lambda = 0$, then $u = 0$ and $w$ can be any nonzero vector in $W$.
	Therefore, the eigenvalues are $\lambda = 1$ and $\lambda = 0$, with corresponding eigenvectors being the nonzero vectors in $U$ for $\lambda = 1$ and the nonzero vectors in $W$ for $\lambda = 0$.
\end{sol}

\begin{exercise}[5a13]
    Suppose $T \in \lin{V}$.
	Suppose $S \in \lin{V}$ is invertible.
    \begin{subexercises}
        \item Prove that $T$ and $S\inv TS$ have the same eigenvalues.
        \item What is the relationship between the eigenvectors of $T$ and the eigenvectors of $S\inv TS$?
    \end{subexercises}
\end{exercise}

\begin{sole}
    \item Let $\lambda$ be an eigenvalue of $T$ with eigenvector $v \in V$, so $T(v) = \lambda v$.
	Then
    \[(S\inv TS)(S\inv v) = S\inv T(v) = S\inv (\lambda v) = \lambda (S\inv v).\]
    Thus, $S\inv v$ is an eigenvector of $S\inv TS$ with eigenvalue $\lambda$.
	This shows that every eigenvalue of $T$ is an eigenvalue of $S\inv TS$.
	Conversely, if $\lambda$ is an eigenvalue of $S\inv TS$ with eigenvector $u$, then $S u$ is an eigenvector of $T$ with eigenvalue $\lambda$.
	Therefore, $T$ and $S\inv TS$ have the same eigenvalues.
    \item The relation between the eigenvectors of $T$ and the eigenvectors of $S\inv TS$ is given by $v \mapsto S\inv v$.
	Specifically, if $v$ is an eigenvector of $T$ with eigenvalue $\lambda$, then $S\inv v$ is an eigenvector of $S\inv TS$ with the same eigenvalue $\lambda$.
	Conversely, if $u$ is an eigenvector of $S\inv TS$ with eigenvalue $\lambda$, then $S u$ is an eigenvector of $T$ with eigenvalue $\lambda$.
\end{sole}

\begin{exercise}[5a14]
    Give an example of an operator on $\rbb^4$ that has no (real) eigenvalues.
\end{exercise}

\begin{sol}
    Consider the linear operator on $\rbb^4$ given by the matrix
    \[T = 
    \begin{pmatrix}
        0 & -1 & 0 & 0 \\
        1 & 0 & 0 & 0 \\
        0 & 0 & 0 & -1 \\
        0 & 0 & 1 & 0
    \end{pmatrix}.
    \]
    Letting $T$ act on $(x_1, x_2, x_3, x_4) \in \rbb^4$, we have
    \[T(x_1, x_2, x_3, x_4) = (-x_2, x_1, -x_4, x_3).\]
    Then any eigenvector $(x_1, x_2, x_3, x_4) \in \rbb^4$ with eigenvalue $\lambda$ must satisfy
    \[(-x_2, x_1, -x_4, x_3) = \lambda (x_1, x_2, x_3, x_4).\]
    This yields $\lambda^2 + 1 = 0$, which has no solution in $\rbb$.
\end{sol}

\begin{exercise}[5a15]
    Suppose $V$ is finite-dimensional, $T \in \lin{V}$, and $\lambda \in \fbb$.
	Show that $\lambda$ is an eigenvalue of $T$ if and only if $\lambda$ is an eigenvalue of the dual operator $T' \in \lin{V'}$.
\end{exercise}

\begin{sol}
    \begin{align*}
        \lambda \text{ is an eigenvalue of } T & \iff T-\lambda I \text{ is not invertible} \\ 
        & \iff (T-\lambda I)' \text{ is not invertible} \\ 
        & \iff T'-\lambda I \text{ is not invertible} \\ 
        & \iff \lambda \text{ is an eigenvalue of } T' \qh
    \end{align*}
\end{sol}

\begin{exercise}[5a16]
    Suppose $v_1, \ldots, v_n$ is a basis of $V$ and $T \in \lin{V}$.
	Prove that if $\lambda$ is an eigenvalue of $T$, then
    \[|\lambda| \leq n \max\{|\mcal(T)_{j, k}| : 1 \leq j, k \leq n\},\]
    where $\mcal(T)_{j, k}$ denotes the entry in row $j$, column $k$ of the matrix of $T$ with respect to the basis $v_1, \ldots, v_n$.
    
    \note{See \exref{6a19} for a different bound on $|\lambda|$.}
\end{exercise}

\begin{sol}
    Let $\lambda$ be an eigenvalue of $T$ with eigenvector $u \in V$, so $Tu = \lambda u$.
	Let $\mcal(T)$ be the matrix of $T$ with respect to the basis $v_1, \ldots, v_n$.
	Then
    \[\lambda u_j = (\mcal(T)u)_j = \sum_{k=1}^n \mcal(T)_{j, k} u_k \quad \text{for each } j = 1, \ldots, n.\]
    Taking absolute values, we get
    \[|\lambda||u_j| = \left|\sum_{k=1}^n \mcal(T)_{j, k} u_k\right| \leq \sum_{k=1}^n |\mcal(T)_{j, k}| |u_k| \leq n \max\{|\mcal(T)_{j, k}| : 1 \leq j, k \leq n\} \max\{|u_k| : 1 \leq k \leq n\}.\]
    Choosing $j$ such that $|u_j| = \max\{|u_k| : 1 \leq k \leq n\}$, we obtain
    \[|\lambda| \leq n \max\{|\mcal(T)_{j, k}| : 1 \leq j, k \leq n\}. \qh\]
\end{sol}

\begin{exercise}[5a17]
    Suppose $\fbb = \rbb$, $T \in \lin{V}$, and $\lambda \in \rbb$.
	Prove that $\lambda$ is an eigenvalue of $T$ if and only if $\lambda$ is an eigenvalue of the complexification $T_\cbb$.
    
    \note{See \emph{\exref{3b33}} for the definition of $T_\cbb$.}
\end{exercise}

\begin{sol}
    Suppose $\lambda$ is an eigenvalue of $T$ with eigenvector $v \in V$, so $Tv = \lambda v$.
	Then in the complexification, we have $T_\cbb(v + i0) = Tv + iT0 = \lambda v + i0 = \lambda (v + i0)$, showing that $\lambda$ is an eigenvalue of $T_\cbb$ with eigenvector $v + i0$.

    Conversely, suppose $\lambda$ is an eigenvalue of $T_\cbb$ with eigenvector $v + iw \in V_\cbb$, where $v, w \in V$.
	Then
    \[T_\cbb(v + iw) = Tv + iTw = \lambda (v + iw) = \lambda v + i\lambda w.\]
    Equating real and imaginary parts, we get
    \[Tv = \lambda v \quad \text{and} \quad Tw = \lambda w.\]
    Since $\lambda \in \rbb$, if $v \neq 0$, then $v$ is an eigenvector of $T$ corresponding to $\lambda$.
	If $v = 0$, then $w \neq 0$ and $Tw = \lambda w$, so $w$ is an eigenvector of $T$ corresponding to $\lambda$.
	This shows that $\lambda$ is an eigenvalue of $T$.
\end{sol}

\begin{exercise}[5a18]
    Suppose $\fbb = \rbb$, $T \in \lin{V}$, and $\lambda \in \cbb$.
	Prove that $\lambda$ is an eigenvalue of the complexification $T_\cbb$ if and only if $\bar{\lambda}$ is an eigenvalue of $T_\cbb$.
\end{exercise}

\begin{sol}
    Suppose $\lambda$ is an eigenvalue of $T_\cbb$ with eigenvector $v + iw \in V_\cbb$, where $v, w \in V$.
	Then
    \[T_\cbb(v + iw) = Tv + iTw = \lambda (v + iw) = \lambda v + i\lambda w.\]
    Moreover, we have
    \[T_\cbb(v - iw) = Tv - iTw.\]
    Note that $v + iw = 0$ if and only if both $v = w = 0$.
	Then $v + iw \neq 0$ implies $v -i w \neq 0$.
	Let $\lambda = a + bi$ with $a, b \in \rbb$.
	Then equating real and imaginary parts, we have
    \[Tv = av - bw \quad \text{and} \quad Tw = bv + aw.\]
    Then
    \begin{align*}
        T_\cbb(v - iw) &= Tv - iTw \\
        &= (av - bw) - i(bv + aw) \\
        &= av - ibv - bw - iaw \\
        &= (a - ib)(v - iw) \\
        &= \bar{\lambda} (v - iw).
    \end{align*}
    Since $v - iw \neq 0$, this shows that $\bar{\lambda}$ is an eigenvalue of $T_\cbb$ with eigenvector $v - iw$.

    Conversely, if $\bar{\lambda}$ is an eigenvalue of $T_\cbb$, then by the same argument, $\lambda$ is an eigenvalue of $T_\cbb$.
	This completes the proof.
\end{sol}

\begin{exercise}[5a19]
    Show that the forward shift operator $T \in \lin{\fbb^\infty}$ defined by
    \[T(z_1, z_2, \ldots) = (0, z_1, z_2, \ldots)\]
    has no eigenvalues.
\end{exercise}

\begin{sol}
    Suppose $\lambda \in \fbb$ is an eigenvalue of $T$ with eigenvector $(z_1, z_2, \ldots) \in \fbb^\infty$.
	Then
    \[T(z_1, z_2, \ldots) = (0, z_1, z_2, \ldots) = \lambda (z_1, z_2, \ldots).\]
    Equating the first component, we get $0 = \lambda z_1$, so either $\lambda = 0$ or $z_1 = 0$.
	If $\lambda = 0$, then the second component gives $z_1 = 0$, and the third component gives $z_2 = 0$, and so on.
	By induction, all $z_k = 0$, contradicting that the eigenvector is nonzero.
	If $\lambda \neq 0$, then $z_1 = 0$, and the second component gives $z_2 = 0$, and so on.
	Again, all $z_k = 0$, a contradiction.
	Therefore, $T$ has no eigenvalues.
\end{sol}

\begin{exercise}[5a20]
    Define the backward shift operator $S \in \lin{\fbb^\infty}$ by
    \[S(z_1, z_2, z_3, \ldots) = (z_2, z_3, \ldots).\]
    \begin{subexercises}
        \item Show that every element of $\fbb$ is an eigenvalue of $S$.
        \item Find all eigenvectors of $S$.
    \end{subexercises}
\end{exercise}

\begin{sole}
    \item Let $\alpha \in \fbb$.
	Consider the vector $(1, \alpha, \alpha^2, \ldots) \in \fbb^\infty$.
	Then
    \[S(1, \alpha, \alpha^2, \ldots) = (\alpha, \alpha^2, \alpha^3, \ldots) = \alpha (1, \alpha, \alpha^2, \ldots),\]
    showing that $\alpha$ is an eigenvalue of $S$ with eigenvector $(1, \alpha, \alpha^2, \ldots)$.
	Since $\alpha \in \fbb$ was arbitrary, every element of $\fbb$ is an eigenvalue of $S$.
    \item Let $\alpha \in \fbb$ and consider the eigenvalue $\alpha$ of $S$.
	Any eigenvector corresponding to $\alpha$ must satisfy
    \[S(z_1, z_2, z_3, \ldots) = \alpha (z_1, z_2, z_3, \ldots),\]
    which gives
    \[(z_2, z_3, z_4, \ldots) = (\alpha z_1, \alpha z_2, \alpha z_3, \ldots).\]
    Equating components, we get $z_2 = \alpha z_1$, $z_3 = \alpha z_2 = \alpha^2 z_1$, and in general $z_k = \alpha^{k-1} z_1$ for $k \geq 1$.
	Therefore, all eigenvectors corresponding to $\alpha$ are of the form $(z_1, \alpha z_1, \alpha^2 z_1, \ldots) = z_1 (1, \alpha, \alpha^2, \ldots)$ with $z_1 \neq 0$.
\end{sole}

\begin{exercise}[5a21]
    Suppose $T \in \lin{V}$ is invertible.
    \begin{subexercises}
        \item Suppose $\lambda \in \fbb$ with $\lambda \neq 0$.
	Prove that $\lambda$ is an eigenvalue of $T$ if and only if $\frac{1}{\lambda}$ is an eigenvalue of $T\inv$.
        \item Prove that $T$ and $T\inv$ have the same eigenvectors.
    \end{subexercises}
\end{exercise}

\begin{sole}
    \item Suppose $\lambda \in \fbb$ with $\lambda \neq 0$ is an eigenvalue of $T$ with eigenvector $v \in V$.
	Then
    \[Tv = \lambda v.\]
    Since $T$ is invertible, we can apply $T\inv$ to both sides to get
    \[v = T\inv (\lambda v) = \lambda T\inv v.\]
    Dividing both sides by $\lambda$ (which is nonzero), we obtain
    \[T\inv v = \frac{1}{\lambda} v.\]
    This shows that $\frac{1}{\lambda}$ is an eigenvalue of $T\inv$ with the same eigenvector $v$.
    
    Conversely, suppose $\mu \in \fbb$ with $\mu \neq 0$ is an eigenvalue of $T\inv$ with eigenvector $v \in V$.
	Then
    \[T\inv v = \mu v.\]
    Applying $T$ to both sides gives
    \[v = \mu Tv \implies Tv = \frac{1}{\mu} v.\]
    This shows that $\frac{1}{\mu}$ is an eigenvalue of $T$ with the same eigenvector $v$.
    
    Therefore, $\lambda$ is an eigenvalue of $T$ if and only if $\frac{1}{\lambda}$ is an eigenvalue of $T\inv$.
    \item Part (a) shows that if $\lambda$ is an eigenvalue of $T$, then $\frac{1}{\lambda}$ is an eigenvalue of $T\inv$ with the same eigenvector.
	Conversely, if $\mu$ is an eigenvalue of $T\inv$, then $\frac{1}{\mu}$ is an eigenvalue of $T$ with the same eigenvector.
	Therefore, $T$ and $T\inv$ have the same eigenvectors.
\end{sole}

\begin{exercise}[5a22]
    Suppose $T \in \lin{V}$ and there exist nonzero vectors $u$ and $w$ in $V$ such that
    \[Tu = 3w \quad \text{and} \quad Tw = 3u.\]
    Prove that $3$ or $-3$ is an eigenvalue of $T$.
\end{exercise}

\begin{sol}
    Consider the vectors $u + w$ and $u - w$.
	Since $u$ and $w$ are both nonzero, at least one of $u + w$ or $u - w$ is nonzero.
	We have
    \[T(u + w) = Tu + Tw = 3w + 3u = 3(u + w),\]
    and
    \[T(u - w) = Tu - Tw = 3w - 3u = -3(u - w).\]
    Hence, one of $u + w$ or $u - w$ is an eigenvector with eigenvalue 3 or $-3$, respectively.
\end{sol}

\begin{exercise}[5a23]
    Suppose $V$ is finite-dimensional and $S, T \in \lin{V}$.
	Prove that $ST$ and $TS$ have the same eigenvalues.
\end{exercise}

\begin{sol}
    Let $\lambda$ be an eigenvalue of $ST$ with eigenvector $v \in V$, so that
    \[STv = \lambda v.\]
    If $\lambda = 0$, then $ST$ is not invertible. Since $V$ is finite-dimensional, $TS$ is also not invertible, so $0$ is an eigenvalue of $TS$.
    Otherwise, if $\lambda \neq 0$, let $w = Tv$.
    Then
    \[Sw = S(Tv) = STv = \lambda v \neq 0 \implies w \neq 0.\]
    Now, we have
    \[TSw = T(STv) = T(\lambda v) = \lambda Tv = \lambda w.\]
    Therefore, $w$ is an eigenvector of $TS$ with eigenvalue $\lambda$.
    Then every eigenvalue of $ST$ is also an eigenvalue of $TS$.
    By symmetry, the same argument shows that if $\lambda$ is an eigenvalue of $TS$, then it is also an eigenvalue of $ST$.
    Hence, $ST$ and $TS$ have the same eigenvalues.
\end{sol}

\begin{exercise}[5a24]
    Suppose $A$ is an $n$-by-$n$ matrix with entries in $\fbb$.
	Define $T \in \lin{\fbb^n}$ by $Tx = Ax$, where elements of $\fbb^n$ are thought of as $n$-by-$1$ column vectors.
    \begin{subexercises}
        \item Suppose the sum of the entries in each row of $A$ equals $1$.
	Prove that $1$ is an eigenvalue of $T$.
        \item Suppose the sum of the entries in each column of $A$ equals $1$.
	Prove that $1$ is an eigenvalue of $T$.
    \end{subexercises}
\end{exercise}

\begin{sole}
    \item Suppose the sum of the entries in each row of $A$ equals $1$.
	Let $v = (1, 1, \ldots, 1)^t \in \fbb^n$.
	Then
    \[Tv = Av = v,\]
    showing that $v$ is an eigenvector of $T$ corresponding to the eigenvalue $1$.
    \item Suppose the sum of the entries in each column of $A$ equals $1$.
	By part (a), this is equivalent to saying that the sum of the entries in each row of $A^t$ equals $1$, hence $1$ is an eigenvalue of the linear operator defined by $A^t$.
	But the eigenvalues of $A^t$ are the same as the eigenvalues of $A$ because $A - \lambda I$ is not invertible if and only if $A^t - \lambda I$ is not invertible.
	Therefore, $1$ is an eigenvalue of $T$.
\end{sole}

\begin{exercise}[5a25]
    Suppose $T \in \lin{V}$ and $u, w$ are eigenvectors of $T$ such that $u + w$ is also an eigenvector of $T$.
	Prove that $u$ and $w$ are eigenvectors of $T$ corresponding to the same eigenvalue.
\end{exercise}

\begin{sol}
    Let $Tu = \lambda u$ and $Tw = \mu w$ for some scalars $\lambda, \mu \in \fbb$.
	Since $u + w$ is also an eigenvector of $T$, we have $T(u + w) = \nu (u + w)$ for some scalar $\nu \in \fbb$.
	Then
    \[\lambda u + \mu w = Tu + Tw = T(u + w) = \nu (u + w) = \nu u + \nu w.\]
    Suppose $u$ and $w$ correspond to different eigenvalues.
    Then the list $u, w$ is linearly independent, and we can equate coefficients of $u$ and $w$ above to obtain $\lambda = \nu = \mu$, which is a contradiction.
    Therefore, $u$ and $w$ must correspond to the same eigenvalue.
\end{sol}

\begin{exercise}[5a26]
    Suppose $T \in \lin{V}$ is such that every nonzero vector in $V$ is an eigenvector of $T$.
	Prove that $T$ is a scalar multiple of the identity operator.
\end{exercise}

\begin{sol}
    Let $v \in V$ be any nonzero vector.
	By assumption, $v$ is an eigenvector of $T$, so $Tv = \lambda_v v$ for some scalar $\lambda_v \in \fbb$.
	Now, let $u \in V$ be another nonzero vector.
	Consider the vector $u + v$.
	If $u$ and $v$ are linearly dependent, then without loss of generality, let $u = \alpha v$ for some scalar $\alpha \in \fbb$.
	Then $Tu = \lambda_u u = \lambda_u \alpha v$, and $Tu = T(\alpha v) = \alpha Tv = \alpha \lambda_v v$.
	Equating these, we get $\lambda_u \alpha v = \alpha \lambda_v v$, and since $\alpha \neq 0$ and $v \neq 0$, it follows that $\lambda_u = \lambda_v$.

    The remaining case is when $u + v$ are linearly independent.
	In this case, $u + v$ is an eigenvector of $T$, so $T(u + v) = \lambda_{u+v} (u + v)$ for some scalar $\lambda_{u+v} \in \fbb$.
	On the other hand, $T(u + v) = Tu + Tv = \lambda_u u + \lambda_v v$.
	Equating these, we get $\lambda_{u+v} u + \lambda_{u+v} v = \lambda_u u + \lambda_v v$.
	Since $u$ and $v$ are linearly independent, we can equate coefficients to get $\lambda_{u+v} = \lambda_u = \lambda_v$.
    
    Since $u$ and $v$ were arbitrary nonzero vectors, it follows that there exists a scalar $\lambda \in \fbb$ such that $Tv = \lambda v$ for all nonzero $v \in V$.
	Hence, $T = \lambda I$.
\end{sol}

\begin{exercise}[5a27]
    Suppose that $V$ is finite-dimensional and $k \in \{1, \ldots, \dim V - 1\}$.
	Suppose $T \in \lin{V}$ is such that every subspace of $V$ of dimension $k$ is invariant under $T$.
	Prove that $T$ is a scalar multiple of the identity operator.
\end{exercise}

\begin{sol}
    Let $v_1 \in V$ be nonzero.
	Our goal is to intersect enough $k$-dimensional subspaces to isolate $\spn v_1$.
	We proceed as follows.
	Let $\dim V = n$, and extend $v_1$ to a basis $v_1, \ldots, v_n$ of $V$.
	For each $i = 2, \ldots, n$, let
    \[U_i = \spn(v_1, w_2, \ldots, w_k), \quad w_2, \ldots, w_k \in \set{v_2, \ldots, v_n} \setminus \set{v_i},\]
    i.e., each $U_i$ must contain $v_1$ as well as $k - 1$ vectors from the basis except $v_i$.
	This choice is possible, since $k \leq n - 1$, so there are enough vectors in $\set{v_2, \ldots, v_n} \setminus \set{v_i}$ to choose $k - 1$ of them.
	The intersection of all such $U_i$ will be exactly $\spn v_1$.
	Since each $U_i$ is invariant under $T$, we have $Tv_1 \in U_i$ for all $i$.
	Therefore, $Tv_1 \in \spn v_1$, which means $Tv_1$ is a scalar multiple of $v_1$.
	Since $v_1$ was arbitrary, it follows that $T$ is a scalar multiple of the identity operator.
\end{sol}

\begin{exercise}[5a28]
    Suppose $V$ is finite-dimensional and $T \in \lin{V}$.
	Prove that $T$ has at most $1 + \dim \range T$ distinct eigenvalues.
\end{exercise}

\begin{sol}
    Let $\lambda_1, \ldots, \lambda_m$ be distinct eigenvalues of $T$ with corresponding eigenvectors $v_1, \ldots, v_m$.
	Since the eigenvalues are distinct, at most one eigenvalue must be 0.
	Without loss of generality, assume $\lambda_1 = 0$ if 0 is an eigenvalue.
    Then $\lambda_2, \ldots, \lambda_m$ are nonzero eigenvalues of $T$.
    For $j = 2, \ldots, m$, we have $Tv_j = \lambda_j v_j$, so $T(\lambda_j\inv v_j) = v_j$, which implies $v_j \in \range T$.
    Therefore, $v_2, \ldots, v_m$ are linearly independent vectors in $\range T$, so $m - 1 \leq \dim \range T$.
    Hence, $m \leq 1 + \dim \range T$.
\end{sol}

\begin{exercise}[5a29]
    Suppose $T \in \lin{\rbb^3}$ and $-4, 5$, and $\sqrt{7}$ are eigenvalues of $T$.
	Prove that there exists $x \in \rbb^3$ such that $Tx - 9x = (-4, 5, \sqrt{7})$.
\end{exercise}

\begin{sol}
    Since $\dim \rbb^3 = 3$, the eigenvalues of $T$ are exactly $-4, 5$, and $\sqrt 7$.
	It follows that 9 is not an eigenvalue of $T$, and $\nul(T - 9I) = \set 0$.
	Then $T - 9I$ is injective, and since $\rbb^3$ is finite-dimensional, $T - 9I$ is also surjective.
	Therefore, for any vector in $\rbb^3$, there exists $x \in \rbb^3$ such that $(T - 9I)x = (-4, 5, \sqrt{7})$.
\end{sol}

\begin{exercise}[5a30]
    Suppose $T \in \lin{V}$ and $(T - 2I)(T - 3I)(T - 4I) = 0$.
	Suppose $\lambda$ is an eigenvalue of $T$.
	Prove that $\lambda = 2$ or $\lambda = 3$ or $\lambda = 4$.
\end{exercise}

\begin{sol}
    Let $v$ be an eigenvector of $T$ corresponding to the eigenvalue $\lambda$.
	Then
    \[(T - 2I)(T - 3I)(T - 4I)v = (T - 2I)(T - 3I)(\lambda - 4)v = (T - 2I)(\lambda - 3)(\lambda - 4)v = (\lambda - 2)(\lambda - 3)(\lambda - 4)v = 0.\]
    Since $v \neq 0$, it must be that $(\lambda - 2)(\lambda - 3)(\lambda - 4) = 0$.
	Therefore, $\lambda = 2$, $\lambda = 3$, or $\lambda = 4$.
\end{sol}

\begin{exercise}[5a31]
    Give an example of $T \in \lin{\rbb^2}$ such that $T^4 = -I$.
\end{exercise}

\begin{sol}
    Consider the rotation matrix by $\pi/4$ radians:
    \[T = 
    \begin{pmatrix}
        \cos(\pi/4) & -\sin(\pi/4) \\
        \sin(\pi/4) & \cos(\pi/4)
    \end{pmatrix}.
    \]
    Then
    \[
    T^4 = 
    \begin{pmatrix}
        \cos(\pi) & -\sin(\pi) \\
        \sin(\pi) & \cos(\pi)
    \end{pmatrix} = 
    \begin{pmatrix}
        -1 & 0 \\
        0 & -1
    \end{pmatrix} = -I. \qh
    \]
\end{sol}

\begin{exercise}[5a32]
    Suppose $T \in \lin{V}$ has no eigenvalues and $T^4 = I$.
	Prove that $T^2 = -I$.
    
    \remark{The proof of this is a little easier if we may use minimal polynomials.}
\end{exercise}

\begin{sol}
    Note that
    \[T^4 - I = (T^2 - I)(T^2 + I) = (T - I)(T + I)(T^2 + I) = 0.\]
    Applying $T^4 - I$ to any nonzero $v \in V$, we get
    \[(T - I)(T + I)(T^2 + I)v = 0.\]
    Let $(T^2 + I)v = w$.
	If $w \neq 0$, then $(T - I)(T + I)w = 0$.
	If $(T + I)w \neq 0$, then
    \[(T - I)(T + I)w = T(T + I)w - I(T + I)w = 0,\]
    which implies $T(T + I)w = (T + I)w$.
	Then $(T + I)w$ would have to be an eigenvector of $T$ corresponding to the eigenvalue $1$, contradicting the assumption that $T$ has no eigenvalues.
	If $(T + I)w = 0$, then $Tw = -w$, which would make $w$ an eigenvector of $T$ corresponding to the eigenvalue $-1$, again contradicting the assumption that $T$ has no eigenvalues.
	Hence, we must have $w = 0$, i.e., $(T^2 + I)v = 0$.
	Since $v$ was arbitrary and nonzero, we conclude that $T^2 + I = 0$, i.e., $T^2 = -I$.
\end{sol}

\begin{exercise}[5a33]
    Suppose $T \in \lin{V}$ and $m$ is a positive integer.
    \begin{subexercises}
        \item Prove that $T$ is injective if and only if $T^m$ is injective.
        \item Prove that $T$ is surjective if and only if $T^m$ is surjective.
    \end{subexercises}
\end{exercise}

\begin{sole}
    \item Suppose $T$ is injective.
	Then if $T^m v = 0$, we have $T(T^{m-1}v) = 0$, which implies $T^{m-1}v = 0$ since $T$ is injective.
	Repeating this argument, we eventually get $v = 0$.
	Hence, $T^m$ is injective.

    Conversely, suppose $T^m$ is injective.
	If $Tv = 0$, then $T^m v = 0$, which implies $v = 0$.
	Hence, $T$ is injective.

    \item Suppose $T$ is surjective.
	Then for any $v \in V$, there exists $w \in V$ such that $Tw = v$.
	Similarly, there exists $w_2 \in V$ such that $Tw_2 = w$, and so on.
	Repeating this $m$ times, we find $u \in V$ such that $T^m u = v$.
	Hence, $T^m$ is surjective.

    Conversely, suppose $T^m$ is surjective.
	Then for any $v \in V$, there exists $u \in V$ such that $T^m u = v$.
	Let $w = T^{m-1} u$.
	Then $Tw = T(T^{m-1} u) = T^m u = v$.
	Hence, $T$ is surjective.
\end{sole}

\begin{exercise}[5a34]
    Suppose $V$ is finite-dimensional and $v_1, \ldots, v_m \in V$.
	Prove that the list $v_1, \ldots, v_m$ is linearly independent if and only if there exists $T \in \lin{V}$ such that $v_1, \ldots, v_m$ are eigenvectors of $T$ corresponding to distinct eigenvalues.
\end{exercise}

\begin{sol}
    Suppose $v_1, \ldots, v_m$ is linearly independent.
	Since $V$ is finite-dimensional, we can extend this list to a basis $v_1, \ldots, v_m, v_{m+1}, \ldots, v_n$ of $V$.
	Define $T \in \lin{V}$ by setting $Tv_i = \lambda_i v_i$ for $i = 1, \ldots, m$, where $\lambda_1, \ldots, \lambda_m$ are distinct scalars, and choose $\lambda_{m+1}, \ldots, \lambda_n$ arbitrarily (distinct from each other and from $\lambda_1, \ldots, \lambda_m$).
	Then $v_1, \ldots, v_m$ are eigenvectors of $T$ corresponding to distinct eigenvalues.

    Conversely, suppose there exists $T \in \lin{V}$ such that $v_1, \ldots, v_m$ are eigenvectors corresponding to distinct eigenvalues $\lambda_1, \ldots, \lambda_m$.
	By 5.11, eigenvectors corresponding to distinct eigenvalues are linearly independent.
	Hence, $v_1, \ldots, v_m$ are linearly independent.
\end{sol}

\begin{exercise}[5a35]
    Suppose that $\lambda_1, \ldots, \lambda_n$ is a list of distinct real numbers.
	Prove that the list $e^{\lambda_1 x}, \ldots, e^{\lambda_n x}$ is linearly independent in the vector space of real-valued functions on $\rbb$.
    
    \hint{Let $V = \spn(e^{\lambda_1 x}, \ldots, e^{\lambda_n x})$, and define an operator $D \in \lin{V}$ by $Df = f'$.
	Find eigenvalues and eigenvectors of $D$.}
\end{exercise}

\begin{sol}
    Suppose there exists a linear combination
    \[a_1 e^{\lambda_1 x} + \cdots + a_n e^{\lambda_n x} = 0\]
    for all $x \in \rbb$.
	Let $V = \spn(e^{\lambda_1 x}, \ldots, e^{\lambda_n x})$ and define $D \in \lin{V}$ by $Df = f'$.
	Then $De^{\lambda_i x} = \lambda_i e^{\lambda_i x}$, so $e^{\lambda_i x}$ is an eigenvector of $D$ corresponding to the eigenvalue $\lambda_i$.
	Since the $\lambda_i$ are distinct, by 5.11, the eigenvectors corresponding to distinct eigenvalues are linearly independent.
	Hence, $a_1 = \cdots = a_n = 0$, proving that $e^{\lambda_1 x}, \ldots, e^{\lambda_n x}$ are linearly independent.
\end{sol}

\begin{exercise}[5a36]
    Suppose that $\lambda_1, \ldots, \lambda_n$ is a list of distinct positive numbers.
	Prove that the list $\cos(\lambda_1 x), \ldots, \\ \cos(\lambda_n x)$ is linearly independent in the vector space of real-valued functions on $\rbb$.
\end{exercise}

\begin{sol}
    Suppose there exists a linear combination
    \[a_1 \cos(\lambda_1 x) + \cdots + a_n \cos(\lambda_n x) = 0\]
    for all $x \in \rbb$.
	Let $V = \spn(\cos(\lambda_1 x), \sin(\lambda_1 x), \ldots, \cos(\lambda_n x), \sin(\lambda_n x))$ and define $D \in \lin{V}$ by $Df = f'$.
	Then $D^2 \cos(\lambda_i x) = -\lambda_i^2 \cos(\lambda_i x)$, so $\cos(\lambda_i x)$ is an eigenvector of $D^2$ corresponding to the eigenvalue $-\lambda_i^2$.
	Since the $\lambda_i$ are distinct and positive, the $-\lambda_i^2$ are distinct, and by 5.11, the eigenvectors corresponding to distinct eigenvalues are linearly independent.
	Hence, $a_1 = \cdots = a_n = 0$, proving that $\cos(\lambda_1 x), \ldots, \cos(\lambda_n x)$ are linearly independent.
\end{sol}

\begin{exercise}[5a37]
    Suppose $V$ is finite-dimensional and $T \in \lin{V}$.
	Define $\mathcal{A} \in \lin{\lin{V}}$ by
    \[\mathcal{A}(S) = TS\]
    for each $S \in \lin{V}$.
	Prove that the set of eigenvalues of $T$ equals the set of eigenvalues of $\mathcal{A}$.
\end{exercise}

\begin{sol}
    Suppose $\lambda$ is an eigenvalue of $T$ with corresponding eigenvector $v \in V$, so $Tv = \lambda v$.
	Consider a nonzero $\phi \in V'$ and consider $S \in \lin{V}$ defined by $Sw = \phi(w)v$ for all $w \in V$.
	Then $S \neq 0$ and for all $w \in V$,
    \[TSw = T(\phi(w)v) = \phi(w)Tv = \phi(w)\lambda v = \lambda \phi(w)v = \lambda Sw.\]
    Hence, $\lambda$ is an eigenvalue of $\mathcal{A}$.
    
    Conversely, if $\lambda$ is an eigenvalue of $\mathcal{A}$ with corresponding eigenvector $S \in \lin{V}$, then $\mathcal{A}(S) = TS = \lambda S$.
	Since $S \neq 0$, there exists $v \in V$ such that $Sv \neq 0$.
	Then $TSv = \lambda Sv$, so $\lambda$ is an eigenvalue of $T$.
	Therefore, the set of eigenvalues of $T$ equals the set of eigenvalues of $\mathcal{A}$.
\end{sol}

\begin{exercise}[5a38]
    Suppose $V$ is finite-dimensional, $T \in \lin{V}$, and $U$ is a subspace of $V$ invariant under $T$.
	The \textit{\emph{\textbf{quotient operator}}} $T / U \in \lin{V / U}$ is defined by
    \[(T / U)(v + U) = Tv + U\]
    for each $v \in V$.
    \begin{subexercises}
        \item Show that the definition of $T / U$ makes sense (which requires using the condition that $U$ is invariant under $T$) and show that $T / U$ is an operator on $V / U$.
        \item Show that each eigenvalue of $T / U$ is an eigenvalue of $T$.
    \end{subexercises}
\end{exercise}

\begin{sole}
    \item  To show that the definition of $T / U$ makes sense, suppose $v + U = v' + U$ for some $v, v' \in V$.
	Then $v - v' \in U$, and since $U$ is invariant under $T$, we have $T(v - v') \in U$, so $Tv - Tv' \in U$, which implies $Tv + U = Tv' + U$.
	Hence, the definition of $T / U$ is well-defined.

    To show that $T / U$ is an operator on $V / U$, note that for any $v + U \in V / U$, $(T / U)(v + U) = Tv + U \in V / U$.
	Linearity follows from the linearity of $T$: for any $\alpha, \beta \in \fbb$ and $v + U, w + U \in V / U$,
    \begin{align*}
        (T / U)(\alpha(v + U) + \beta(w + U)) &= (T / U)(\alpha v + \beta w + U) \\
        &= T(\alpha v + \beta w) + U \\
        &= \alpha Tv + \beta Tw + U \\
        &= \alpha(T / U)(v + U) + \beta(T / U)(w + U)
    \end{align*}
    \item Suppose $\lambda$ is an eigenvalue of $T / U$ with corresponding eigenvector $v + U \in V / U$, so $(T / U)(v + U) = Tv + U = \lambda(v + U) = \lambda v + U$.
	This means $Tv - \lambda v = (T - \lambda I)v \in U$.
    Suppose $\lambda$ is not an eigenvalue of $T$, hence $(T - \lambda I)$ is injective.
    
    Consider the restriction $(T - \lambda I)|_U$, which is injective since $T - \lambda I$ is injective.
    Since $U$ is finite-dimensional, the injectivity of $(T - \lambda I)|_U$ implies that it is also surjective onto $U$.
    This implies that there exists $u \in U$ such that $(T - \lambda I)u = (T - \lambda I)v$, which means $(T - \lambda I)(v - u) = 0$.
    Since $v - u \neq 0$ (otherwise $v \in U$, which would contradict $v + U$ being an eigenvector of $T / U$), this contradicts the injectivity of $T - \lambda I$.
    Hence, $\lambda$ must be an eigenvalue of $T$.
\end{sole}

\begin{exercise}[5a39]
    Suppose $V$ is finite-dimensional and $T \in \lin{V}$.
	Prove that $T$ has an eigenvalue if and only if there exists a subspace of $V$ of dimension $\dim V - 1$ that is invariant under $T$.
\end{exercise}

\begin{sol}
    Suppose $T$ has an eigenvalue $\lambda$ with corresponding eigenvector $v \in V$.
	Consider $\range(T - \lambda I)$.
	Then $\dim \range(T - \lambda I) \leq \dim V - 1$.
	Extend $\range(T - \lambda I)$ to a subspace $U$ of $V$ of dimension $\dim V - 1$.
	Then $U$ is invariant under $T$ because for any $u \in U$, $Tu = (T - \lambda I)u + \lambda u \in U$.
	Hence, there exists a subspace of $V$ of dimension $\dim V - 1$ that is invariant under $T$.

    Conversely, suppose there exists a subspace $U$ of $V$ of dimension $\dim V - 1$ that is invariant under $T$.
	Then $V / U$ is one-dimensional, and the induced operator $T / U$ on $V / U$ must have an eigenvalue $\lambda$.
	By \exref{5a38} (b), any eigenvalue of $T / U$ is also an eigenvalue of $T$.
	Hence, $T$ has an eigenvalue.
\end{sol}

\begin{exercise}[5a40]
    Suppose $S, T \in \lin{V}$ and $S$ is invertible.
	Suppose $p \in \pcal(\fbb)$ is a polynomial.
	Prove that
    \[p(STS\inv) = Sp(T)S\inv.\]
\end{exercise}

\begin{sol}
    We first prove by induction that $(STS\inv)^k = ST^kS\inv$ for all non-negative integers $k$.
	For $k = 0$, we have $(STS\inv)^0 = I = ST^0S\inv$.
	Suppose the result holds for some $k = n$, so that $(STS\inv)^n = ST^nS\inv$.
	Then
    \begin{align*}
        (STS\inv)^{n+1} &= (STS\inv)^n (STS\inv) \\
        &= (ST^nS\inv)(STS\inv) \\
        &= S T^n (S\inv S) T S\inv \\
        &= S T^n T S\inv \\
        &= S T^{n+1} S\inv.
    \end{align*}
    By induction, the result holds for all non-negative integers $k$.

    Now let $p(z) = a_0 + a_1 z + \cdots + a_m z^m$ be any polynomial.
	Then
    \begin{align*}
        p(STS\inv) &= a_0 I + a_1 (STS\inv) + \cdots + a_m (STS\inv)^m \\
        &= a_0 I + a_1 (STS\inv) + \cdots + a_m (ST^mS\inv) \\
        &= S(a_0 I + a_1 T + \cdots + a_m T^m)S\inv \\
        &= Sp(T)S\inv.
    \end{align*}
    Hence, the result holds for all polynomials $p$.
\end{sol}

\begin{exercise}[5a41]
    Suppose $T \in \lin{V}$ and $U$ is a subspace of $V$ invariant under $T$.
	Prove that $U$ is invariant under $p(T)$ for every polynomial $p \in \pcal(\fbb)$.
\end{exercise}

\begin{sol}
    Let $u \in U$.
	Since $U$ is invariant under $T$, we have $T(u) \in U$.
	By induction, suppose $T^k(u) \in U$ for some non-negative integer $k$.
	Then $T^{k+1}(u) = T(T^k(u)) \in U$ since $U$ is invariant under $T$.
	Hence, $T^k(u) \in U$ for all non-negative integers $k$.
	Now let $p(z) = a_0 + a_1 z + \cdots + a_m z^m$ be any polynomial.
	Then
    \[p(T)(u) = a_0 u + a_1 T(u) + \cdots + a_m T^m(u) \in U,\]
    since each term is in $U$.
	Therefore, $U$ is invariant under $p(T)$ for every polynomial $p \in \pcal(\fbb)$.
\end{sol}

\begin{exercise}[5a42]
    Define $T \in \lin{\fbb^n}$ by $T(x_1, x_2, x_3, \ldots, x_n) = (x_1, 2x_2, 3x_3, \ldots, nx_n)$.
    \begin{subexercises}
        \item Find all eigenvalues and eigenvectors of $T$.
        \item Find all subspaces of $\fbb^n$ that are invariant under $T$.
    \end{subexercises}
    \apnote{5a42}{ for a generalization to any diagonal operator.
	The solution here remains as to show the student the mechanism.}
\end{exercise}

\begin{sole}
    \item Let $\lambda$ be an eigenvalue of $T$ with corresponding eigenvector $v = (x_1, x_2, \ldots, x_n) \in \fbb^n$.
	Then $T(v) = \lambda v$, which means
    \[(x_1, 2x_2, 3x_3, \ldots, nx_n) = (\lambda x_1, \lambda x_2, \ldots, \lambda x_n).\]
    Comparing components, we get $\lambda x_i = i x_i$ for each $i = 1, 2, \ldots, n$.
	Hence, if $x_i \neq 0$, then $\lambda = i$.
	Therefore, the eigenvalues of $T$ are $1, 2, \ldots, n$, and the corresponding eigenvectors are nonzero multiples of the standard basis vectors $e_1, e_2, \ldots, e_n$.
    \item Let $U$ be an invariant subspace of $\fbb^n$ under $T$.
	Let
    \[u = \sum_{i = 1}^n \alpha_i e_i \in U,\]
    where $e_i$ are the standard basis vectors of $\fbb^n$ and $\alpha_i \in \fbb$.
	Since $U$ is invariant under $T$, we have
    \[Tu = \sum_{i = 1}^n \alpha_i T(e_i) = \sum_{i = 1}^n \alpha_i i e_i \in U.\]
    Consider some index $i$ such that $\alpha_i \neq 0$.
	Then
    \[u_1 = Tu - i u = \sum_{j = 1}^n \alpha_j j e_j - i \sum_{j = 1}^n \alpha_j e_j = \sum_{j = 1}^n \alpha_j (j - i) e_j \in U.\]
    The $i$-th term in this sum is $\alpha_i (i - i) e_i = 0$, hence the $e_i$ term is eliminated from $u_1$.
	Consider another index $j$ of $u_1$ such that $\alpha_j \neq 0$.
	Then we can form
    \[u_2 = Tu_1 - j u_1 = \sum_{k = 1}^n \alpha_k (k - i) k e_k - j \sum_{k = 1}^n \alpha_k (k - i) e_k = \sum_{k = 1}^n \alpha_k (k - i) (k - j) e_k \in U.\]
    The $j$-th term in this sum is $\alpha_j (j - i) (j - j) e_j = 0$, and $u_2$ now has all terms except the $i$-th and $j$-th terms.
	By repeating this process for each index with a nonzero coefficient, we can isolate each $\alpha_i e_i$ in $U$ whenever $\alpha_i \neq 0$.
    This shows that for any $u \in U$, all the standard basis vectors $e_i$ corresponding to nonzero coefficients in $u$ must also be in $U$.
	Therefore, $U$ must be the span of some subset of the standard basis vectors $\{e_1, e_2, \ldots, e_n\}$.
\end{sole}

\begin{exercise}[5a43]
    Suppose that $V$ is finite-dimensional, $\dim V > 1$, and $T \in \lin{V}$.
	Prove that $\set{p(T) \mid p \in \pcal(\fbb)} \neq \lin{V}$.
\end{exercise}

\begin{sol}
	Let $v_1, \ldots, v_n$ be a basis for $V$, with $n = \dim V$.
	Consider the operators $T_1, T_2 \in \lin V$ defined by
    \[T_1v_1 = v_2, \quad T_1v_i = 0~\text{for}~i = 2, \ldots, n,\]
    and
    \[T_2v_1 = v_1, \quad T_2v_i = 0~\text{for}~i = 2, \ldots, n.\]
    Then $T_2T_1v_1 = T_2(T_1v_1) = T_2v_2 = 0$, and $T_1T_2v_1 = T_1(T_2v_1) = T_1v_1 = v_2$.
	Hence, $T_2T_1 \neq T_1T_2$, so $T_1$ and $T_2$ do not commute.
	Let $S = \set{p(T) \mid p \in \pcal(\fbb)}$. Let $p, q \in \pcal$ be defined as
    \[p(z) = a_0 + a_1 z + \dots + a_m z^m, \quad q(z) = b_0 + b_1 z + \dots + b_k z^k.\]
    Since $p(T)q(T) = q(T)p(T)$, it follows that all operators in $S$ commute.
    However, if $T_1 = r(T)$ and $T_2 = s(T)$ for some polynomials $r, s \in \pcal$, then $T_1$ and $T_2$ would commute, which is a contradiction.
    Hence, $T_1$ and $T_2$ cannot both be in $S$, so $S \neq \lin{V}$.
\end{sol}